% beautybook — decorative book class with coloured theorem boxes and a
% two-tone chapter-heading treatment, adapted from the CTAN/GitHub
% BeautyLaTeX/latex-template package.
%
% Simplified from upstream: the upstream ships four cover styles (cn/en/
% enfig/birkar) and a 36MB set of photographic chapter-corner decorations
% (inner_pics/titleimages/, Pixabay-style stock photos with no separate
% licence grant of their own). This starter uses ONLY the text-only "en"
% cover (no background photo) and a small self-drawn abstract mark
% (inner_pics/beautybook-mark.png, original artwork made for this starter)
% in place of every upstream photo, so nothing here carries an unverified
% third-party image licence.
\documentclass[lang=en,12pt]{beautybook}
\RequirePackage[utf8]{inputenc}
\RequirePackage{times}
\RequirePackage[T1]{fontenc}
\RequirePackage[brazil]{babel}
\RequirePackage{microtype}
\RequirePackage[round]{natbib}

\definecolor{coverbgcolor}{HTML}{e0e0e0}
\definecolor{coverfgcolor}{HTML}{1f3134}
\definecolor{coverbar}{HTML}{7c9092}
\definecolor{bottomcolor}{HTML}{2c4f54}
\definecolor{nuanbai}{HTML}{f5f5f5}
\coverstyle={
    cover-choose=en,
}
\mathstyle={
    math-font=plain,
}

\mynewtheorem{
    defi={\textbf{Definição}}[section]{interior style={left color=ReD!8,right color=ReD!5!CyaN!50}, borderline west={1.5mm}{0mm}{ReD}},
    thm={\textbf{Teorema}}[section]{interior style={left color=CyaN!80!black!20,right color=CyaN!80!black!15!CyaN!50}, borderline west={1.5mm}{0mm}{CyaN!80!black}},
    lem={\textbf{Lema}}[section]{interior style={left color=BluE!8,right color=BluE!5!CyaN!50}, borderline west={1.5mm}{0mm}{BluE}},
    prop={\textbf{Proposição}}[section]{interior style={left color=OrangE!8,right color=OrangE!5!CyaN!50}, borderline west={1.5mm}{0mm}{OrangE}},
    exam={\textbf{Exemplo}}[chapter]{interior style={left color=DarkGreen!8,right color=DarkGreen!5!CyaN!50}, borderline west={1.5mm}{0mm}{DarkGreen}},
}
\newtheorem*{remark}{\textbf{Observação}}

\begin{document}
\thispagestyle{empty}
\title{Curso Maker}
\subtitle{Módulo 2}
\edition{}
\bookseries{}
\author{Poggers Maker}
\pressname{beautybook}
\presslogo{inner_pics/beautybook-mark.png}
\makecover

\makeatletter
\definecolor{bg}{HTML}{e0e0e0}
\definecolor{fg}{HTML}{2c4f54}
\colorlet{outermarginbgcolor}{bg}
\colorlet{outermarginfgcolor}{fg}
\oddoutermargin{\sffamily \leftmark}
\evenoutermargin{\sffamily\@title}
\chapimage{inner_pics/beautybook-mark.png}
\makeatother

\colorlet{framegolden}{fg}
\colorlet{framegray}{bg!50}

\frontmatter
\pagenumbering{Roman}

{\thispagestyle{empty}
\chapter*{Prefácio}
This short book collects the definitions and first theorems of graph
colouring, aimed at a reader who has seen basic graph theory but not yet
chromatic numbers.

\hfill
\begin{tabular}{lr}
    &-- Poggers Maker\\
    & 2026
\end{tabular}
\clearpage}

\thispagestyle{empty}
\tableofcontents\let\cleardoublepage\clearpage

\mainmatter
\pagenumbering{arabic}

\partabstract{\hspace*{2em} The first part develops the basic vocabulary
of graph colouring and proves the elementary upper bounds every later
result builds on.}
\part{Foundations}

\chapter{Proper Colourings}
\label{ch:colourings}

\section{Definitions}

\begin{defi}[Proper colouring]
\label{def:colouring}
A \emph{proper $k$-colouring} of a graph $G = (V, E)$ is a function
$c \colon V \to \{1, \dots, k\}$ such that $c(u) \neq c(v)$ whenever
$uv \in E$.
\end{defi}

\begin{defi}[Chromatic number]
\label{def:chromatic}
The \emph{chromatic number} $\chi(G)$ is the least $k$ for which $G$
admits a proper $k$-colouring.
\end{defi}

\begin{thm}[Greedy bound]
\label{thm:greedy-bound}
For any graph $G$ with maximum degree $\Delta$, $\chi(G) \leq \Delta + 1$.
\end{thm}

\begin{proof}
Order the vertices arbitrarily and colour them in that order, at each step
assigning the least colour not used by an already-coloured neighbour.
Each vertex has at most $\Delta$ neighbours, so at most $\Delta$ colours
are forbidden, leaving one of the first $\Delta + 1$ colours free.
\end{proof}

\begin{exam}[An odd cycle]
\label{exam:cycle}
The cycle $C_5$ has $\Delta = 2$, so Theorem~\ref{thm:greedy-bound} gives
$\chi(C_5) \leq 3$. In fact $\chi(C_5) = 3$: two colours are never enough
for a cycle of odd length, since the colours would have to alternate and
cannot close up consistently.
\end{exam}

\begin{lem}[Bipartite graphs]
\label{lem:bipartite}
A graph $G$ is bipartite if and only if $\chi(G) \leq 2$.
\end{lem}

\begin{remark}
Lemma~\ref{lem:bipartite} makes $2$-colourability and bipartiteness the
same question, and both can be checked in linear time by breadth-first
search.
\end{remark}

\chapter{The Four Colour Theorem}
\label{ch:four-colour}

\begin{prop}[Planar upper bound]
\label{prop:planar-bound}
Every planar graph satisfies $\chi(G) \leq 4$.
\end{prop}

Proposition~\ref{prop:planar-bound} is the Four Colour Theorem, first
proved by computer-assisted case analysis
by~\citet{appel1977four}, and later given a more tractable
proof by~\citet{robertson1997four}.

\bibliographystyle{plainnat}
\bibliography{references}

\end{document}
